\documentclass[journal]{IEEEtran}

\usepackage{booktabs}
\usepackage{float}
\usepackage{url}
\usepackage[hidelinks]{hyperref}

\begin{document}

\title{Supplementary Material for ``ContractKV: Certified Prompt Compilation
for Transactional Prefix Reuse in Multi-Tenant LLM Services''}

\author{Bin~Zhou and Xi~Cao,~\IEEEmembership{Member,~IEEE}}

\maketitle

\section{Current-vLLM Relevance Audit}

\subsection{Protocol and Runtime Identity}

This audit tests whether ContractKV's structural opportunity and integration
objects remain present in a current runtime. It does not repeat the full service
matrix or claim that the frozen vLLM~0.10.2 PSTP patch applies unchanged.

We use the official installed \texttt{vllm==0.25.1+cu129} wheel, Qwen3-8B in
FP16, and one NVIDIA RTX~3090. Native automatic prefix caching is enabled. Five
RULER \texttt{niah\_single\_1} cases and five LongBench HotpotQA cases are
replayed under raw order and the same previously qualified frozen layout. Each
cell runs in three fresh server processes, yielding 60 measured requests.
Serving uses \texttt{max\_model\_len=20000},
\texttt{max\_num\_seqs=1}, and a 64-token output limit.

The optional FlashInfer sampler is disabled because the host's default CUDA
toolkit predates the wheel runtime. Both layouts use this setting, all caches
and logs reside on the data volume, and no vLLM source file is modified. The
audited scheduler and cache-manager files match the installed wheel's
\texttt{RECORD} hashes.

\subsection{Real Serving Results}

\begin{table}[H]
\centering
\caption{Native vLLM~0.25.1 APC performance. Token counts come from server metrics.}
\label{tab:vllm025-results}
\scriptsize
\setlength{\tabcolsep}{2.7pt}
\begin{tabular}{llrrrr}
\toprule
Workload & Layout & Mean & p95 & Hit & New \\
 & & TTFT & TTFT & tok. & tok. \\
\midrule
RULER & Raw & 295.66 & 302.48 & 32 & 1,172 \\
RULER & Compiled & \textbf{59.59} & \textbf{66.56} &
1,136 & \textbf{68} \\
HotpotQA & Raw & 4,754.28 & 5,402.30 & 32 & 15,658.2 \\
HotpotQA & Compiled & \textbf{153.60} & \textbf{174.79} &
15,616 & \textbf{74.2} \\
\bottomrule
\end{tabular}
\end{table}

\begin{table}[H]
\centering
\caption{Paired behavior audit.}
\label{tab:vllm025-quality}
\small
\begin{tabular}{llrr}
\toprule
Workload & Layout & Task pass & Pair equal \\
\midrule
RULER & Raw & 100\% & -- \\
RULER & Compiled & 100\% & 100\% \\
HotpotQA & Raw & 60\% & -- \\
HotpotQA & Compiled & 60\% & 100\% \\
\bottomrule
\end{tabular}
\end{table}

The compiled layout reduces mean TTFT by 79.84\% on RULER and 96.77\% on
HotpotQA. Observed new-prefill work falls by 94.20\% and 99.53\%,
respectively. All requests complete. Every paired raw/compiled output is
identical, and both the mean and minimum paired task-score differences are
zero.

The absolute HotpotQA task pass rate is 60\% in both cells. It is neither an
official LongBench aggregate nor a ContractKV-induced regression. This small
audit uses task scores only to test paired behavior preservation; the main
paper reports the full quality protocols.

\subsection{Source Portability Audit}

\begin{table}[H]
\centering
\caption{vLLM~0.25.1 integration-object audit.}
\label{tab:vllm025-hooks}
\scriptsize
\setlength{\tabcolsep}{3pt}
\begin{tabular}{p{0.23\columnwidth}p{0.64\columnwidth}}
\toprule
Requirement & Status \\
\midrule
Layout admission & Available (\texttt{Scheduler.add\_request}) \\
Prefix observation & Available (\texttt{get\_computed\_blocks}) \\
Slot allocation & Available (\texttt{allocate\_slots}) \\
Cache publication & Available; interface adaptation required \\
Partial-block registration & Available; remains prefix-chain metadata \\
Frozen patch compatibility & Not drop-in \\
\bottomrule
\end{tabular}
\end{table}

Native vLLM still searches a longest prefix chain rather than arbitrary
non-contiguous stable segments. The admission, lookup, allocation, and
publication object classes remain available, but their signatures and
publication internals differ from the frozen implementation. The compiler's
frozen output can therefore be consumed by the current runtime, whereas PSTP
requires a version-specific backend adaptation.

This audit supports architectural relevance and adaptation feasibility, not
full implementation equivalence across versions. Machine-readable request
records, native metric summaries, source hashes, and the protocol manifest are
available in the public artifact:
\url{https://github.com/zhouz4884-coder/contractkv}.

\end{document}
